% =====================================================================
% Cinco problemas adicionales para «Problemas Cuantitativos Propuestos»
% del capítulo 3. Van dentro del mismo \begin{enumerate} ... \end{enumerate},
% justo después del problema del cometa (momento angular), y continúan la
% numeración (11 a 15).
% Requiere: amsmath, siunitx (ya cargados en main.tex).
% =====================================================================

    \item \textbf{Impulso y ley de la fuerza:} Una pelota de tenis de
    \SI{57}{\gram} llega horizontalmente a una pared con una rapidez de
    \SI{30}{\meter\per\second} y rebota en la misma dirección, en sentido
    contrario, con \SI{40}{\meter\per\second}. El contacto con la pared dura
    \SI{5.0}{\milli\second}.
    \begin{enumerate}
        \item Tomando como positivo el sentido de salida, calcule el cambio en
        la cantidad de movimiento $\Delta p$ de la pelota.
        \item Utilizando la ley de la fuerza, $F = \Delta p/\Delta t$, calcule
        la magnitud de la fuerza media que la pared ejerce sobre la pelota y
        compárela con el peso de la pelota.
        \textit{[Resp: a) $\Delta p = \SI{3.99}{\kilo\gram\meter\per\second}$;
        b) $F \approx \SI{8.0e2}{\newton}$, unas $1400$ veces el peso]}
    \end{enumerate}

    \item \textbf{Trabajo con la Fuerza Promedio:} Se quiere subir un satélite
    de \SI{500}{\kilo\gram} desde la superficie terrestre hasta una órbita a
    \SI{400}{\kilo\meter} de altura. (Datos: $M_T = \SI{5.97e24}{\kilo\gram}$,
    $R_T = \SI{6371}{\kilo\meter}$,
    $G \approx \SI{6.674e-11}{\newton\meter\squared\per\kilo\gram\squared}$).
    \begin{enumerate}
        \item Calcule la fuerza promedio que debe vencerse en el ascenso,
        usando el promedio geométrico de las fuerzas gravitacionales en los
        extremos, $\langle F\rangle = \sqrt{F(r_a)\,F(r_b)} = GM_Tm/(r_a r_b)$.
        \item Calcule el trabajo $W = \langle F\rangle\,(r_b - r_a)$ y
        verifique que coincide con el cambio de energía potencial
        $\Delta U = GM_Tm\,(1/r_a - 1/r_b)$.
        \item Compare con la aproximación $mgh$, con
        $g = \SI{9.8}{\meter\per\second\squared}$. ¿En qué porcentaje la
        sobreestima? \textit{[Resp: a) $\langle F\rangle \approx
        \SI{4.62e3}{\newton}$; b) $W \approx \SI{1.85e9}{\joule}$;
        c) $mgh = \SI{1.96e9}{\joule}$, un \SI{6}{\percent} mayor, del orden
        de $h/R_T$]}
    \end{enumerate}

    \item \textbf{Energía Potencial Gravitacional y Velocidad de Escape:}
    Utilizando la conservación de la energía mecánica con
    $U(r) = -GMm/r$ y la condición de que el cuerpo llegue al infinito con
    rapidez nula ($K + U = 0$), calcule la velocidad de escape desde la
    superficie de la Luna. (Datos: $M_L = \SI{7.35e22}{\kilo\gram}$,
    $R_L = \SI{1737}{\kilo\meter}$). Compárela con la de la Tierra,
    $\SI{11.2}{\kilo\meter\per\second}$, y explique por qué la Luna no ha
    podido retener una atmósfera. \textit{[Resp:
    $v_{\text{esc}} = \sqrt{2GM_L/R_L} \approx \SI{2.38}{\kilo\meter\per\second}$,
    unas cinco veces menor que la de la Tierra]}

    \item \textbf{Tercera Ley de Kepler y Centro de Masas:} Júpiter
    ($m_J = \SI{1.90e27}{\kilogram}$) orbita al Sol
    ($M_\odot = \SI{1.99e30}{\kilo\gram}$, $R_\odot = \SI{6.96e8}{\meter}$) a
    una distancia media $a = \SI{5.20}{UA} = \SI{7.78e11}{\meter}$.
    \begin{enumerate}
        \item Usando la tercera ley de Kepler en la forma
        $T^{2} \propto a^{3}$ y tomando como referencia la órbita terrestre
        ($a = \SI{1}{UA}$, $T = 1$~año), calcule el período orbital de
        Júpiter en años.
        \item Con la definición del centro de masas, calcule la distancia
        entre el centro del Sol y el centro de masas del sistema Sol--Júpiter
        y exprésela en radios solares. ¿Está dentro o fuera del Sol?
        \item Calcule la rapidez con la que el Sol gira alrededor de ese
        centro de masas. (Este bamboleo del Sol es el que permite descubrir
        planetas en otras estrellas por el método de la velocidad radial).
        \textit{[Resp: a) $T \approx \num{11.9}$~años; b)
        $r_\odot = m_J a/(M_\odot + m_J) \approx \SI{7.42e8}{\meter} \approx
        1.07\,R_\odot$, fuera del Sol; c) $v_\odot = 2\pi r_\odot/T \approx
        \SI{12.5}{\meter\per\second}$]}
    \end{enumerate}

    \item \textbf{Momento Angular y Precesión:} Una rueda de bicicleta de
    masa $m = \SI{2.0}{\kilo\gram}$ y radio $R = \SI{0.25}{\meter}$, con
    toda su masa en el aro ($I = mR^{2}$), se cuelga de una cuerda por un
    extremo de su eje, a una distancia $d = \SI{0.36}{\meter}$ del centro de
    la rueda, con el eje horizontal. La rueda gira a
    $\num{5.0}$~vueltas por segundo.
    \begin{enumerate}
        \item Calcule el momento angular $L = I\omega$ de la rueda y el torque
        $\tau = mgd$ que el peso ejerce respecto al punto de suspensión.
        \item Como el torque es perpendicular a $\boldsymbol{L}$, la relación
        $\Delta\boldsymbol{L} = \boldsymbol{\tau}\,\Delta t$ hace girar el eje
        en un plano horizontal en lugar de tumbarlo. Calcule la rapidez
        angular de esta precesión, $\Omega = \tau/L$, y el tiempo que tarda el
        eje en dar una vuelta completa.
        \item ¿Qué pasa con ese tiempo si la rueda gira el doble de rápido?
        ¿Y si deja de girar? \textit{[Resp: a)
        $L \approx \SI{3.93}{\kilo\gram\meter\squared\per\second}$,
        $\tau \approx \SI{7.06}{\newton\meter}$; b)
        $\Omega \approx \SI{1.80}{\radian\per\second}$, una vuelta cada
        \SI{3.5}{\second}; c) se duplica; sin giro, $L = 0$ y la rueda
        simplemente cae]}
    \end{enumerate}
